%% kapa: exam=1 round=32 year=2021 period=1 form=A subject=economics subject_ko=경제학원론 seq=16 number=56
%% source: Q-Net 감정평가사 기출문제 (cst00309) artlSeq=5208397 file=extracted/1008/5208397/A형/2021년도 제32회 감정평가사 1차 1교시 A형.pdf
%% conversion: kordoc 4.24.51 -> tools/gen1.py (auto); review=proofread-sonnet; flags=none
\begin{question}{56}
독점기업의 가격차별 전략 중 이부가격제(two-part pricing)에 관한 설명으로 옳은 것을 모두 고른 것은?
\begin{framedtext}
ㄱ. 서비스 요금 설정에서 기본요금(가입비)과 초과사용량 요금(사용료)을 분리하여 부과하는 경우가 해당된다. \par
ㄴ. 적은 수량을 소비하는 소비자의 평균지불가격이 낮아진다. \par
ㄷ. 소비자잉여는 독점기업이 부과할 수 있는 가입비의 한도액이다. \par
ㄹ. 자연독점 하의 기업이 평균비용 가격설정으로 인한 손실을 보전하기 위해 선택한다. \par
\end{framedtext}
\begin{choices}
  \item ㄱ, ㄴ
  \item ㄱ, ㄷ
  \item ㄴ, ㄷ
  \item ㄱ, ㄴ, ㄷ
  \item ㄴ, ㄷ, ㄹ
\end{choices}
\end{question}
\answer{2}
